\section{Hodge theory and formal lifting}\label{proj:sec:hodge}

We retain the geometric package of
Proposition~\ref{proj:prop:signed-construction}.  In particular,
\(\mathcal P=\sqrt[\ell]{\OO_{\PP^s}(1)}\) is a root gerbe, \(C\) is
its tautological line bundle, and
\[
 g:S\longrightarrow\mathcal Y
 \xrightarrow{p}\mathcal P
\]
is a representable projective morphism followed by a smooth projective
bundle.  The Stein image \(\mathcal T\subset\mathcal Y\) is finite over
\(\mathcal P\), and \(g_*\OO_S=\OO_{\mathcal T}\).  The reduced Cartier
divisor \(S\subset\mathcal X\) has dimension \(N=n-1\).  If
\(J=(S\cap T)_{\mathrm{red}}\), then \(S,J\) are Du Bois,
\(\mathcal I_{J/S}\) is maximal Cohen--Macaulay, and
\begin{equation}\label{proj:hg:geometric-lines}
 \mathcal B:=
 \mathcal Hom_S(\mathcal I_{J/S},\omega_S)
 \simeq \OO_S(aS)=g^*C^a,
 \qquad
 \OO_S(S)=g^*C,
\end{equation}
where \(a>0\) is an integer.  Off \(T\), the fixed ambient isomorphism
\(\omega_{\mathcal X}(S)\simeq\OO_{\mathcal X}(aS)\) supplies these
identifications by adjunction.  We prove that functions and transverse
parameters can be lifted through every infinitesimal neighborhood of
\(S\).

\subsection{A global vanishing statement}

We use graded-polarizable mixed Hodge modules on ordinary complex
algebraic varieties.  The needed established results are projective
strictness and the usual functorial operations
\cite[Theorem~2.14 and Section~4]{Proj-Saito1990}, the comparison with the
Du Bois complex and its coherent dual
\cite[Theorem~0.2 and Corollary~0.3]{Proj-Saito2000}, and
Kodaira--Saito vanishing \cite[Proposition~2.33]{Proj-Saito1990}.
No Hodge-module theory on stacks is assumed.  Our objects on smooth
stacks will be compatible systems on scheme etale charts, whose
filtered differential modules descend.

Throughout this section, differential modules are right modules with
increasing filtration.  The module itself is in degree zero in its
Spencer de Rham complex.  Thus, on a smooth chart \(V\), the term in
degree \(-i\) of \(\Gr^F_k\DR_V(M)\) is
\[
 \Gr^F_{k-i}M\otimes_{\OO_V}\bigwedge^iT_V.
\]
In particular, if \(F_{<0}M=0\), the lowest graded de Rham complex is
the sheaf \(F_0M\) in degree zero.

\begin{lemma}\label{proj:hg:negative-vanishing}
Let \(M\) be a compatible system of mixed Hodge modules in perverse
degree zero on the scheme etale charts of \(\mathcal Y\).  Suppose
its support is finite over \(\mathcal P\).  Then, for every integer
\(k\), every positive integer \(b\), and every \(j<0\),
\[
 \mathbb H^j\!\left(
   \mathcal Y,\Gr^F_k\DR_{\mathcal Y}(M)\otimes p^*C^{-b}
 \right)=0.
\]
Here the graded de Rham complex is the descended coherent complex.
\end{lemma}

\begin{proof}
Because \(p\) is finite on the support, its direct image has only
perverse cohomology in degree zero.  Denote that system on
\(\mathcal P\) by \(M'\).  Projective strictness on the ordinary
scheme charts gives
\begin{equation}\label{proj:hg:graded-direct-image}
 Rp_*\Gr^F_k\DR_{\mathcal Y}(M)
 \simeq \Gr^F_k\DR_{\mathcal P}(M').
\end{equation}
We first explain why this is a global identity of descended complexes,
rather than a claim of local vanishing.

For right differential modules, the direct image is formed from the
filtered transfer bimodule
\[
 \mathcal D_{\mathcal Y\to\mathcal P}
 =
 \OO_{\mathcal Y}\otimes_{p^{-1}\OO_{\mathcal P}}
                         p^{-1}\mathcal D_{\mathcal P},
\]
with its order filtration, by derived tensor over
\(\mathcal D_{\mathcal Y}\) and derived direct image.  These operations
may equivalently be performed on Rees modules.  Applying de Rham on
the base means further derived tensor with \(\OO_{\mathcal P}\),
whose filtration begins in degree zero.  The filtered Spencer
resolution, associativity of derived tensor, and projection formula
identify the result with the direct image of de Rham upstairs:
the transfer module tensored over the base differential operators
with its structure sheaf is \(\OO_{\mathcal Y}\).
These are canonical sheaf constructions on the etale site; the
Spencer resolutions are locally free over differential operators.
They therefore respect chart changes before passing to associated
gradeds.  Strictness and the concentration in perverse degree zero
can be checked on the base charts by the ordinary projective
direct-image theorem.  Coherent cohomology on those charts is
unchanged by using their etale sites.  This proves
Equation~\eqref{proj:hg:graded-direct-image} as an identity that computes
global coherent hypercohomology.  The same argument applies to the
representable finite morphism used next.

There is a representable finite surjective morphism
\[
 v:\PP^s_z\longrightarrow\mathcal P
\]
defined by the power map
\([z_0:\cdots:z_s]\mapsto[z_0^\ell:\cdots:z_s^\ell]\)
and the root \(\OO_{\PP^s_z}(1)\).  In particular,
\(v^*C=\OO_{\PP^s_z}(1)\).  On the standard affine opens of \(\PP^s\)
choose the trivial root, obtaining scheme etale charts
\(\check U_i\to\mathcal P\).  On \(z_i\ne0\), the morphism \(v\)
factors through \(\check U_i\).  Pull back \(M'\) on these charts
and take perverse cohomology in degree zero.  Functoriality and the
chosen root identify these objects on overlaps, giving a system
\(H\) on the Zariski opens of \(\PP^s_z\).

For completeness, this is a global graded-polarizable mixed Hodge
module on \(\PP^s_z\).  The perverse, filtered, and weight data glue.
The pure weight gradeds have strict-support decompositions, which
glue by uniqueness.  On a smooth connected dense stratum of an
irreducible support, a polarization from a nonempty Zariski open
extends as a flat pairing: the fundamental group of that open
surjects onto the fundamental group of the stratum.  Its Hodge
compatibility extends by continuity; nondegeneracy and positivity
persist for the extended flat pairing.  Quasi-unipotence at the
boundary is already supplied by the local Hodge modules.  Saito's
extension theorem for polarizable variations
\cite[Section~3.b]{Proj-Saito1990} and uniqueness of strict-support
extension identify this global pure object with the glued one.
There is no additional boundary at infinity, since \(\PP^s_z\) is
projective.  The weight extensions then give the stated mixed
object.

The system \(M'\) is a retract of \(v_+H\).  To see this on a chart,
use the \emph{whole} base-changed finite cover.  Over \(\check U_i\)
it is a disjoint union of copies of the corresponding affine
coordinate chart of \(\PP^s_z\).  The unit on unshifted constants
and its dual trace, using smooth duality in equal dimensions, have
composition multiplication by the degree.  On the finite etale
locus the trace is summation over sheets.  On each connected
smooth base chart, the shifted constant Hodge module is the
rank-one intersection complex and has endomorphism ring \(\Q\);
its endomorphisms are determined on a dense open.  Thus the
composition equals the degree globally, with no additional term
supported on the branch locus.  Ramification is therefore
allowed.  Tensoring with \(M'\), applying proper
projection formula, and taking perverse cohomology in degree zero
gives the claimed retraction, since finite direct image is
perverse exact.  Divide the trace by the degree.  All maps are
canonical under etale base change, so the retraction descends also
on filtered differential modules.

Now apply the finite version of
Equation~\eqref{proj:hg:graded-direct-image} and projection formula.
The desired hypercohomology is a direct summand of
\[
 \mathbb H^j\!\left(
   \PP^s_z,\Gr^F_k\DR(H)\otimes\OO_{\PP^s_z}(-b)
 \right).
\]
This is zero for \(j<0\) by ordinary negative-ample
Kodaira--Saito vanishing.  One may apply it to the pure weight
gradeds and then use the exact sequences of the weight filtration;
the Hodge filtrations of those sequences are strict.
\end{proof}

\subsection{The lowest Hodge piece}

Take a separated quasi-projective scheme etale chart
\(U\to\mathcal Y\), of finite type.  Write
\(S_U=S\times_{\mathcal Y}U\), \(J_U=J\times_{\mathcal Y}U\),
and again \(g:S_U\to U\) for the induced projective morphism.
For \(j:S_U\setminus J_U\hookrightarrow S_U\), put
\[
 A^\bullet_U=\mathbf D\!\left(
             j_!\Q^H_{S_U\setminus J_U}[N]\right),
 \qquad
 A_{0,U}={}^{p}\!H^0A^\bullet_U,
 \qquad
 M_U={}^{p}\!H^1g_*A_{0,U}.
\]
All these constructions commute with etale restriction.  The \(M_U\)
form a system supported on \(\mathcal T\), to which
Lemma~\ref{proj:hg:negative-vanishing} applies.

\begin{lemma}\label{proj:hg:lowest-piece}
The following identifications hold compatibly on the charts:
\[
 F_{<0}A_{0,U}=0,\qquad F_0A_{0,U}=\mathcal B|_{S_U},
 \qquad
 F_{<0}M_U=0,\qquad
 F_0M_U=R^1g_*(\mathcal B|_{S_U}).
\]
In a smooth ambient embedding \(S_U\hookrightarrow V\) of
codimension \(c\), the underlying module of \(A_{0,U}\) is
\(\mathcal H^c_{S_U}(\omega_V)\), localized off \(J_U\).
The lowest-piece inclusion is the adjunction, or Ext-to-support,
inclusion off \(J_U\), extended by meromorphic localization.
\end{lemma}

\begin{proof}
The difference triangle for the constant objects of \(S_U\) and
\(J_U\), together with Du Bois comparison, identifies its
degree-zero graded de Rham complex with
\(\mathcal I_{J_U/S_U}\).  Coherent duality and the
maximal-Cohen--Macaulay property therefore give
\begin{equation}\label{proj:hg:derived-lowest}
 \Gr^F_k\DR(A^\bullet_U)=0\quad(k<0),
 \qquad
 \Gr^F_0\DR(A^\bullet_U)=\mathcal B|_{S_U}[0].
\end{equation}
The shift \([N]\) cancels the dimension shift in the dualizing
complex.  In the smooth case this convention says
\(\mathbf D(\Q^H[N])=\Q^H[N](N)\); its right differential module
has lowest piece \(\omega\) at index zero.

We spell out the passage from the complex to perverse cohomology.
Let \(q\) be the smallest filtration index occurring in any
perverse cohomology object of \(A^\bullet_U\), locally on a fixed
chart.  Such a lower bound exists because the complex is bounded
and its filtrations are good.  At index \(q\), the graded Spencer
complex of every perverse cohomology object has only its
degree-zero term.  The spectral sequence from perverse truncation
therefore has a single nonzero row and gives
\[
 \mathcal H^i\Gr^F_q\DR(A^\bullet_U)
   =F_q\,{}^p\!H^iA^\bullet_U.
\]
If \(q<0\), this contradicts Equation~\eqref{proj:hg:derived-lowest}.
Thus all these objects have \(F_{<0}=0\).  Applying the same
argument at index zero gives \(F_0A_{0,U}=\mathcal B|_{S_U}\),
and gives zero for the lowest piece of the other perverse
cohomology objects.  No perversity assertion about the shifted
constant complex is required.

For the unfiltered description, duality identifies the dual of the
reduced constant complex, in the ambient \(V\), with
\[
 R\Gamma_{[S_U]}(\omega_V)[c].
\]
Its degree-zero module is the first nonzero support cohomology
\(\mathcal H^c_{S_U}(\omega_V)\).  The support of \(J_U\) is locally
set-theoretically principal in \(S_U\), by the construction.
Invert a local defining function, lifted to \(V\); this exact
meromorphic localization realizes \(j_*\) on the underlying
module.  It yields the asserted description of \(A_{0,U}\).

We also need compatibility of the two descriptions of its lowest
piece.  On the smooth locus of \(S_U\setminus J_U\), filtered
duality and smooth graph pushforward identify it with the ordinary
dualizing sheaf included by residue.  This is the Ext-to-support
inclusion, with the normalization fixed by the ambient adjunction
isomorphism.  Naturality in smooth etale coordinates fixes the
same scalar on every chart.  Agreement extends over
\(S_U\setminus J_U\): the first support-cohomology module has no
sections supported on a smaller-dimensional subset.  For example,
this follows from the Cousin resolution of the smooth dualizing
sheaf, whose first term supported on \(S_U\) is a sum of injective
modules at its codimension-\(c\) generic points.  It extends across
\(J_U\) by localization.  Hence the inclusion is the one stated
in the lemma, not merely an abstract isomorphism of coherent
sheaves.

Finally use a graph embedding followed by a smooth projection to
compute \(g_+\).  At filtration zero the relative Spencer complex
has only its top term, namely the direct image of
\(\mathcal B|_{S_U}\) on the graph.  Projective strictness gives
\[
 F_{<0}M_U=0,\qquad
 F_0M_U=R^1g_*(\mathcal B|_{S_U}),
\]
and identifies this sheaf with its actual submodule in the
unfiltered degree-one direct image.  The constructions used here
are canonical on etale changes of charts.
\end{proof}

\subsection{Lifting through the infinitesimal neighborhoods}

Let \(I=\OO_{\mathcal X}(-S)\).  The etale object \(S_U\to S\)
extends uniquely and compatibly over every nilpotent thickening
\(qS\); denote its extension by \(qS_U\).
These are quasi-projective schemes by
Proposition~\ref{proj:prop:signed-construction}: they are separated
and quasi-finite over \(X\), hence are open in schemes finite over
the projective variety \(X\).  Powers and quotients of \(I\) below
are pulled back to the respective thickenings.

\begin{proposition}\label{proj:prop:formal-lifting}
For every such chart \(U\), every \(j\geq0\), and every \(k\geq1\),
the transition
\[
 g_*(I^j/I^{j+k+1})
       \longrightarrow g_*(I^j/I^{j+k})
\]
is surjective as a map of sheaves of abelian groups.  Its kernel is
\(\OO_{\mathcal T_U}\otimes C^{-j-k}\), where
\(\mathcal T_U=\mathcal T\times_{\mathcal Y}U\).
The statements are compatible with further etale changes of
charts.  On an affine chart, these transitions are also
surjective on global sections.
\end{proposition}

\begin{proof}
Define
\[
 E_{j,k}=g_*(I^j/I^{j+k}),\qquad j\geq0,\quad k\geq1.
\]
Here \(g_*\) uses the underlying topological map.  We do not
assume that a thickening already has a ringed-space map to \(U\).
The first graded layer is
\[
 E_{j,1}=\OO_{\mathcal T_U}\otimes C^{-j}.
\]
The exact sequence for adjacent lengths has kernel
\(g_*\OO_{S_U}\otimes C^{-j-k}\), and its connecting homomorphism
takes values in
\[
 R^1g_*\OO_{S_U}\otimes C^{-j-k}.
\]

Induct on \(k\), proving the assertion simultaneously for all \(j\).
Assume all shorter transitions are surjective.  A section of
\(E_{j,k}\) with zero length-one term comes from \(E_{j+1,k-1}\)
when \(k>1\).  By induction it locally lifts from \(E_{j+1,k}\),
so its connecting class is zero.  Also \(E_{j,k}\to E_{j,1}\)
is surjective.  For \(k=1\) the factorization is immediate.
The obstruction therefore factors through a map \(D_k\) on the
graded algebra of length-one layers, with graded shift \(k\).
Choose local lifts and take their Cech differences.  The
difference of a product is the sum of the two first-order
differences; products of the errors vanish in the layer at issue.
Thus \(D_k\) is a \(\C\)-derivation from this graded algebra to
the graded \(R^1g_*\OO_{S_U}\)-module.

In degree zero, restrict it along
\(\OO_U\to\OO_{\mathcal T_U}\).  The resulting derivation factors
through \(\Omega_U^1\), hence defines a section of
\(T_U\otimes R^1g_*\OO_{S_U}\otimes C^{-k}\).
These local sections descend to
\begin{equation}\label{proj:hg:obstruction-section}
 e\in H^0\!\left(
 \mathcal Y,T_{\mathcal Y}\otimes R^1g_*\OO_S\otimes C^{-k}
 \right)
 =
 H^0\!\left(
 \mathcal Y,T_{\mathcal Y}\otimes F_0M\otimes C^{-(a+k)}
 \right).
\end{equation}
The equality uses Equation~\eqref{proj:hg:geometric-lines},
Lemma~\ref{proj:hg:lowest-piece}, and projection formula.
For descent, lifts of functions pull back on the unique etale
thickenings.  Their Cech classes pull back by ordinary flat base
change for the coherent obstruction sheaves on the reductions.
Differentials pull back and generate under an etale map, proving
the asserted compatibility.

Shrink \(U\) and take etale coordinates \(t_1,\ldots,t_d\),
where \(d=\dim U\), and a frame of \(C^{-1}\).  Its pullback is
a conormal frame, denoted by \(y\).  By induction lift the
coordinates along \(g\) modulo \(I^k\), and lift \(y\) to
\(I/I^{k+1}\).  On an open cover of \(S_U\), lift these one step
further, obtaining functions \(h_{i,\lambda}\) modulo \(I^{k+1}\)
and generators \(y_i\) modulo \(I^{k+2}\).  Formal etaleness of
the coordinate map makes each \(h_i\) a local map from
\((k+1)S_U\) to \(U\).  On overlaps write
\begin{equation}\label{proj:hg:coordinate-errors}
 h_{j,\lambda}-h_{i,\lambda}
       =e_{ij,\lambda}y_i^k,
 \qquad
 y_j-y_i=\eta_{ij}y_i^{k+1}.
\end{equation}
The reduced coefficients are Cech cocycles representing
\(D_k(t_\lambda)\) and \(D_k(y)\) in the chosen frames.
Let \(\sigma=y^{-a}\) denote the induced frame of
\(\mathcal B|_{S_U}\).

Off \(J_U\), the fixed adjunction isomorphism gives
\[
 \omega_{(k+1)S_U}
 \simeq\OO_{\mathcal X}((a+k)S)|_{(k+1)S_U}.
\]
It therefore defines local dualizing sections
\(b_i=y_i^{-(a+k)}\).  Under the trace inclusion
\(\omega_{S_U}\hookrightarrow\omega_{(k+1)S_U}\), the identities
\begin{equation}\label{proj:hg:dualizing-errors}
 b_j-b_i=-(a+k)\eta_{ij}\sigma,\qquad
 y_i^kb_i=\sigma,\qquad y_i^{k+1}b_i=0
\end{equation}
hold when the dualizing sheaves are realized in support
cohomology.  Indeed the first identity is the binomial expansion
of \((1+\eta_{ij}y_i^k)^{-(a+k)}\); its quadratic terms vanish
because \(2k\geq k+1\).  The other two identities describe
Cartier trace and the annihilator of the thickening.

Embed \((k+1)S_U\) as a closed subscheme of a smooth open
\(Z\subset\overline Z=\PP^m\).  The graph of the reduced map
\(g\) is closed in \(\overline Z\times U\), by its properness over
\(U\).  Use this reduced graph for an ambient realization of
\(A_{0,U}\).  Each local graph lift \(h_i\) maps \(b_i\), by
Ext-to-support cohomology, to a section \(\delta_i(b_i)\) of the
underlying graph module.  We claim that
\begin{equation}\label{proj:hg:graph-identity}
 \delta_j(b_j)-\delta_i(b_i)
 =
 -(a+k)\eta_{ij}\sigma
 +\sum_{\lambda=1}^d
          (e_{ij,\lambda}\sigma)\cdot\partial_{t_\lambda}.
\end{equation}
The reduced dualizing sections on the right use the inclusion
of Lemma~\ref{proj:hg:lowest-piece}.

Here is the local calculation, including the non-smooth case.
Put \(c_0=\dim Z-N\).
The Ext-to-support map identifies the dualizing sheaf of the
thickening with its annihilator submodule in
\(\mathcal H^{c_0}_{S_U}(\omega_Z)\).  One can obtain this from
the support-cohomology spectral sequence: there is no support
cohomology below \(c_0\), so in total degree \(c_0\) the Ext
group is the homomorphisms from the thickening's structure sheaf
into that first support-cohomology module.  Thus the calculation
does not require \(S_U\) to be a local complete intersection in
\(Z\).

Lift \(h_{i,\lambda}\) locally to functions on \(Z\).  In
\(Z\times U\), the graph inclusion on a thickness-dualizing
section \(b\) is the generalized fraction
\[
 \frac{b\,dt_1\wedge\cdots\wedge dt_d}
      {\prod_{\lambda=1}^d(t_\lambda-h_{i,\lambda})}.
\]
This is successive support cohomology in the additional
coordinate equations, or equivalently the Koszul residue for
the graph.  It may be computed etale-locally near the graph
branch.  After first taking support cohomology from \(Z\), the
translated new coordinates form a regular sequence, so their
support cohomology is concentrated in their top number.
Localization in either set of translated coordinates gives the
same localization on this support-torsion module: the two
translations differ nilpotently on every section.

Changing from \(h_i\) to \(h_j\) therefore gives a finite Taylor
expansion on \(b_j\).  Products of two differences in
Equation~\eqref{proj:hg:coordinate-errors} annihilate \(b_j\), again
because \(2k\geq k+1\).  Its linear term uses
\[
 (h_{j,\lambda}-h_{i,\lambda})b_j
       =e_{ij,\lambda}\sigma.
\]
For a top differential form the right action satisfies
\[
 \left(\frac{dt_\lambda}{t_\lambda-h}\right)
       \cdot\partial_{t_\lambda}
   =\frac{dt_\lambda}{(t_\lambda-h)^2}.
\]
Consequently the doubled pole has the positive sign displayed in
Equation~\eqref{proj:hg:graph-identity}.  Together with
Equation~\eqref{proj:hg:dualizing-errors}, this proves that identity.
The calculation initially takes place off \(J_U\).  The sections
and identities extend meromorphically across \(J_U\) in the
\emph{unfiltered} localized module; arbitrary finite pole orders
there are allowed.

Now push by \(\overline Z\times U\to U\), using relative
Spencer.  If \(i\) denotes the closed reduced graph embedding and
\(\pi\) the projection, closed direct image is perverse exact and
\(\pi_+i_+A_{0,U}=g_+A_{0,U}\).  Thus degree-one holonomic
differential-module cohomology of this Spencer direct image is
precisely the underlying module of
\(M_U={}^{p}\!H^1g_*A_{0,U}\).
The cochain \((\delta_i(b_i))\) lies in its top,
degree-zero term, which has no outgoing relative Spencer
differential.  Its Cech difference is therefore a boundary in
total degree one.  A sufficiently refined ambient cover around
the closed graph computes this assertion; all sheaves involved
are supported there.  The reduced cocycles in
Equation~\eqref{proj:hg:graph-identity} represent their
\(R^1g_*\mathcal B\) classes inside \(F_0M\), by
Lemma~\ref{proj:hg:lowest-piece}.  Right differentiation in the
\(U\)-coordinates acts on the pushforward complex.
Hence the right side of Equation~\eqref{proj:hg:graph-identity}
represents zero in the unfiltered module \(M\).

That right side lies in \(F_1M\); the term involving \(\eta\)
lies in \(F_0M\).  Since \(F_1M\) is an actual subsheaf of \(M\),
its vanishing in \(M\) is vanishing in \(F_1M\).  Taking the
order-one symbol thus says that the global section \(e\) of
Equation~\eqref{proj:hg:obstruction-section} is killed by
\begin{equation}\label{proj:hg:symbol-map}
 T_{\mathcal Y}\otimes F_0M
           \longrightarrow \Gr^F_1M,
\end{equation}
after twisting by \(C^{-(a+k)}\).
This reasoning places no Hodge-filtration bound on the
auxiliary cochain \((b_i)\).

There is no term preceding degree \(-1\) in
\(\Gr^F_1\DR(M)\), because \(F_{<0}M=0\).  Its degree-\(-1\)
hypercohomology after the indicated twist is therefore precisely
the space of global sections of the kernel of
Equation~\eqref{proj:hg:symbol-map}.  Since \(a+k>0\),
Lemma~\ref{proj:hg:negative-vanishing} makes that space zero.
We conclude that \(e=0\).

In a local coordinate and line trivialization, this says that
every class \([e_{ij,\lambda}\sigma]\) is already zero in
\(F_0M\), hence in \(M\).  Its actual differential-operator image
is consequently zero, before taking symbols.  The unfiltered
relation now gives
\[
 (a+k)[\eta_{ij}\sigma]=0.
\]
Lowest-piece injectivity and \(a+k>0\) imply \(D_k(y)=0\).
The degree-zero derivation is zero as well: it vanishes on
\(\OO_U\), which surjects onto \(\OO_{\mathcal T_U}\).
Locally the whole graded length-one algebra is generated by
that coefficient algebra and the conormal frame \(y\).
Thus \(D_k=0\) in every degree, proving the induction and all
sheaf transitions.

Their kernels are the coherent sheaves
\(\OO_{\mathcal T_U}\otimes C^{-j-k}\).  On an affine \(U\)
these have no first cohomology.  The exact sequences of sheaves
of abelian groups therefore also give surjectivity on global
sections.  Choosing lifts successively gives compatible formal
lifts of any chosen functions on \(\mathcal T_U\), and, when
\(C\) is trivialized, of its transverse conormal frame.
\end{proof}
